\section{Introduction}
\label{sec:intro}

Law is a domain in which the quality of an answer depends not only on whether it sounds
reasonable, but also on whether it is supported by the correct legal authority. A useful legal
information system must identify relevant statutes and precedents, interpret them in context,
account for their relationships, and present conclusions that can be checked against
authoritative sources. These requirements make legal artificial intelligence (LegalAI) a
distinctive research challenge. Unlike many general question-answering tasks, legal information
retrieval can depend on exact statutory language, the relationship between a provision and its
subordinate clauses, the authority of a judicial decision, and the jurisdiction and date to
which a legal rule applies. A response that overlooks one of these conditions may appear
convincing while remaining legally incomplete or incorrect.

The effort to automate legal reasoning predates modern machine learning. Early research explored
how computational systems could represent legal rules, compare cases, and construct arguments
resembling those used by lawyers. Ashley's work on case-based reasoning, including HYPO,
demonstrated how a system could reason about legal disputes by comparing fact patterns and
identifying legally meaningful similarities and differences \citep{ashley1991}. This tradition
established an important principle: legal reasoning cannot always be reduced to matching a
question with a sentence. It often requires identifying which facts matter, which rules apply,
and why one authority supports an argument more strongly than another. However, rule-based and
case-based approaches also faced practical constraints. Their performance depended on how legal
knowledge was represented, how exceptions and changing authorities were handled, and how much
domain-specific knowledge could be encoded and maintained.

The growth of statistical natural language processing and pretrained language models changed how
researchers approached these problems. Rather than relying exclusively on manually encoded rules,
machine-learning systems began to learn linguistic patterns from legal documents and perform
tasks such as document classification, information extraction, and case retrieval. Domain-adapted
models such as LEGAL-BERT illustrated the value of training language representations on legal text
rather than relying only on general-purpose language data \citep{chalkidis2020}. Benchmarks such
as LexGLUE and LegalBench subsequently provided more systematic ways to evaluate legal language
models across different tasks and forms of legal reasoning \citep{chalkidis2022, guha2023}. These
developments helped move LegalAI from narrowly engineered expert systems toward more flexible,
data-driven models.

The arrival of large language models (LLMs) expanded these capabilities further. Modern models can
summarize legal documents, explain provisions in natural language, compare arguments, and draft
responses to complex questions. Yet fluency introduces a serious reliability problem. A model may
generate an answer that is coherent and professionally written without accurately reflecting the
law, may cite a nonexistent authority, or may combine correct legal statements in a way that
produces an unsupported conclusion. Research on legal hallucinations has shown why apparent
linguistic quality cannot be treated as evidence of legal correctness \citep{dahl2024,
magesh2025}. In legal practice, this distinction is particularly important because users need to
verify not only the answer, but also the authorities and reasoning on which it depends.

Retrieval-augmented generation (RAG) offers one response to this challenge. Instead of asking an
LLM to answer exclusively from its learned parameters, a RAG system retrieves external documents
and supplies their content as evidence for answer generation. This design can make legal answers
more current, traceable, and specific to an authoritative corpus. However, retrieval does not
automatically guarantee correctness. If a system retrieves an irrelevant judgment, misses a
decisive statutory exception, or separates a provision from its legal hierarchy, the generator
may produce a plausible answer from incomplete evidence. Likewise, the presence of a citation
does not establish that the cited passage actually supports the associated claim. Reliable legal
RAG therefore depends on the interaction between retrieval quality, document representation,
evidence selection, generation, and verification, not on the language model alone.

A central challenge is that legal documents are structured sources of authority, not merely
collections of text passages. Statutes contain parts, chapters, sections, articles, clauses,
provisos, and exceptions; judgments contain facts, issues, arguments, reasoning, and decisions.
Flattening these documents into independent text chunks can remove the context needed to
interpret a retrieved passage correctly. A short provision may appear relevant when read in
isolation but mean something different when its parent section, exception, amendment, or related
provision is considered. Structure-aware indexing attempts to preserve this context by attaching
information about a passage's position in the original document. Hybrid retrieval, which combines
lexical matching with semantic similarity, addresses a complementary problem: exact legal
expressions and section numbers can be important, while conceptually relevant evidence may use
different wording from the user's question. Rank-fusion methods can combine these complementary
retrieval signals, although their effectiveness must be evaluated rather than assumed.

Recent research reflects the growing breadth of this field. Broad surveys now examine LegalAI
through the roles of data, modelling, retrieval, generation, and evaluation, while more focused
studies investigate case retrieval, legal information extraction, and the reliability of legal
language models \citep{feng2024, premasiri2025, kelsall2025, chi2026}. These reviews provide
essential foundations, but the field still benefits from a connected account of how legal AI's
design priorities have evolved: from representing legal reasoning, to learning legal language, to
retrieving external evidence, and finally to preserving the structure and provenance needed to
verify generated answers.

This survey follows that evolution. Rather than treating every model or benchmark as an isolated
development, it examines the relationship between legal representation, retrieval architecture,
evidence grounding, and reliability. It considers classical case-based reasoning, legal language
models and benchmarks, legal RAG, structure-aware and hybrid retrieval, hallucination and
citation verification, and the challenges created by jurisdictional and linguistic variation.
Particular attention is given to systems that must operate over structured legal authorities and
to jurisdictions that are less represented in dominant English-language datasets. Pakistani legal
information retrieval provides a useful motivating context because its authoritative sources
include constitutional provisions, legislation, procedural laws, and judicial decisions whose
relationships must be preserved during retrieval. It is considered here as a specific instance of
a wider challenge: a legal AI system cannot be considered reliable merely because its language
model is powerful if its evidence representation fails to preserve the legal context in which that
evidence must be interpreted.

The survey makes four contributions:
\begin{enumerate}[leftmargin=1.4em, itemsep=2pt, topsep=2pt]
\item \textbf{Historical synthesis:} It connects early computational legal reasoning with
domain-specific language models and contemporary retrieval-augmented systems, explaining the
problems that motivated each transition.
\item \textbf{Architectural analysis:} It distinguishes document representation, retrieval, rank
fusion, generation, and evidence verification, clarifying how these components contribute to the
overall reliability of legal AI.
\item \textbf{Critical comparison:} It examines contemporary legal AI systems in terms of their
design choices, evaluation methods, reported strengths, and limitations, while distinguishing
results obtained under different datasets and experimental conditions.
\item \textbf{Research agenda:} It identifies open challenges in structure-preserving retrieval,
jurisdiction-aware evidence selection, multilingual legal access, knowledge-graph integration, and
verifiable answer generation.
\end{enumerate}

These contributions are organized around the research questions in Table~\ref{tab:rq}.

\begin{table*}[t]
\centering
\caption{Research questions guiding the survey.}
\label{tab:rq}
\footnotesize
\setlength{\tabcolsep}{5pt}
\renewcommand{\arraystretch}{1.3}
\begin{tabularx}{\textwidth}{>{\raggedright\arraybackslash}p{4.0cm} >{\raggedright\arraybackslash}X >{\raggedright\arraybackslash}X}
\toprule
\textbf{Research question} & \textbf{Focus} & \textbf{Purpose} \\
\midrule
\textbf{RQ1.} How has legal AI evolved? & From case-based reasoning and rule-based systems to legal language models and RAG & Identify the principal methodological transitions and the limitations that motivated them. \\
\textbf{RQ2.} How do legal retrieval architectures affect evidence quality? & Lexical, dense, hybrid, rank-fusion, structure-aware, and knowledge-graph-based retrieval & Explain how systems locate relevant authorities and preserve the context needed to interpret them. \\
\textbf{RQ3.} How should reliability be assessed in legal AI? & Hallucination, factual support, citation correctness, retrieval quality, and human evaluation & Distinguish fluent generation from evidence-supported legal answers and examine the limits of common evaluation practices. \\
\textbf{RQ4.} What challenges remain for jurisdiction-aware and multilingual legal AI? & Legal hierarchy, jurisdictional variation, English and low-resource languages, and evidence provenance & Identify research opportunities for systems that must retrieve and verify legal information across different legal and linguistic contexts. \\
\bottomrule
\end{tabularx}
\end{table*}

Together, these questions frame the survey around a single problem: how can legal AI move beyond
generating plausible legal language toward producing answers whose supporting evidence is
relevant, contextually correct, and independently verifiable? The next section explains how the
literature was identified, assessed, and organized to investigate this question.
